\section{Utility Conditionals: finite bridges}
\label{sec:annex-p19-rationality-6-4}

\begin{annexitem}{annex:p19.r64.choice-gap}{thm:p19-r64-choice-gap}{Two available acts and the missing alternative}
\begin{proof}
Subtract the two defined finite expectations:
\[
 U_1-U_0=(r_1-r_0)v-(c_1-c_0).
\]
Over the menu consisting exactly of these two acts, a positive difference
has the unique maximizer \(a_1\), a negative difference the unique maximizer
\(a_0\), and a zero difference makes both maximal. This proves the
strict comparison and tie boundary, without adding a rule for choosing
among tied maxima. Replacing \(v\) by \(v+\delta\), with \(\delta>0\),
changes the difference by \((r_1-r_0)\delta\), whose sign gives all
three asserted cases. Replacing \(r_1\) by \(r_1+\varepsilon\) within
\([0,1]\), with \(\varepsilon>0\) and fixed \(v>0\), changes the
difference by \(\varepsilon v>0\).

For underdetermination hold \(r_1=.8\), \(v=10\), \(c_1=3\),
\(c_0=0\) fixed. At \(r_0=.2\), the difference is \(3>0\);
at \(r_0=.9\), it is \(-4<0\). Thus even supplying the cost of
the first act does not eliminate the missing-alternative dependency.
Holding both probabilities instead and changing \(c_1\) can likewise
reverse the comparison. All probabilities lie in the permitted interval.
The argument uses the stipulated additive consequence utilities and
two-act maximization only; it is not a representation theorem for arbitrary
preferences or a psychological prediction without those assumptions.
\end{proof}
\end{annexitem}

\begin{annexitem}{annex:p19.r64.diagnostic-failure}{thm:p19-r64-diagnostic-failure}{Association, intervention, and the choice value}
\begin{proof}
There are two equally probable background units. When \(Z=0\), the
equations give \((D,Q)=(0,0)\); when \(Z=1\), they give
\((D,Q)=(1,1)\). Hence \(P(D=1)=1/2>0\), and ordinary ratio
conditioning gives \(P(Q=1\mid D=1)=(1/2)/(1/2)=1\).
Under the finite intervention convention, setting \(D=d\) replaces only
the equation for \(D\) by the constant \(d\). The equation \(Q=Z\)
and the background law are unchanged, so \(P(Q=1\mid\operatorname{do}(D=d))
=1/2\) for both \(d=0,1\).

The declared consequence utility is the function \(d-10q\).
After setting \(D=1\), its two equally probable values are \(1,-9\),
with expectation \(-4\). After setting \(D=0\), they are \(0,-10\),
with expectation \(-5\). The former is larger. Thus the observational
conditional can be certain while the outcome's intervention probability
is unaffected by the act, and the added direct benefit can make the act
preferable. No conditioning on a null antecedent and no inference of
causal structure from a marginal table alone is used.
\end{proof}
\end{annexitem}
